\section*{Capítulo \ref{ch:g9:alternating-current} --- A corrente alternada}

\begin{solution}{exo:g9:alternating-current:1}
A CC circula sempre num sentido sob uma \omterm{def:g8:voltage-voltmeter:voltage}{tensão} constante --- o regime
da \omterm{def:g3:first-electric-circuit:battery}{pilha}. A CA inverte periodicamente sob uma \omterm{def:g8:voltage-voltmeter:voltage}{tensão} que oscila --- o
regime da tomada.
\end{solution}

\begin{solution}{exo:g9:alternating-current:2}
$T$: a \omterm{def:g2:measuring-time:duration}{duração} de um ciclo completo; $f$: os ciclos por segundo;
$f = 1/T$ e $T = 1/f$. Para \qty{50}{Hz}: $T = \qty{0.02}{s}$. Para
$T = \qty{0.01}{s}$: $f = \qty{100}{Hz}$.
\end{solution}

\begin{solution}{exo:g9:alternating-current:3}
Cem vezes: duas inversões em cada um dos cinquenta ciclos.
\end{solution}

\begin{solution}{exo:g9:alternating-current:4}
$T = 4 \times 5 = \qty{20}{ms} = \qty{0.02}{s}$;
$f = 1/0.02 = \qty{50}{Hz}$.
\end{solution}

\begin{solution}{exo:g9:alternating-current:5}
\qty{230}{V} é a \omterm{prop:g9:alternating-current:effective}{tensão eficaz} --- o valor contínuo constante que
aqueceria igualmente bem, e o que os \omterm{def:g8:voltage-voltmeter:voltmeter}{voltímetros} de CA leem;
\qty{325}{V} é o pico, a crista que a oscilação de fato toca, para
cada lado, cinquenta vezes por segundo.
\end{solution}

\begin{solution}{exo:g9:alternating-current:6}
O aquecimento não liga a mínima para o sentido --- a \omterm{def:g8:ohms-law:resistance}{resistência}
esquenta em qualquer marcha, e por isso as chaleiras bebem a
oscilação diretamente. Carregar é química ao contrário, e a química
exige um sentido consistente: o bloco converte a oscilação em
constância primeiro.
\end{solution}

\begin{solution}{exo:g9:alternating-current:7}
A bateria: uma linha plana em \qty{9}{V}. A fonte alternada: uma onda
suave oscilando entre $+9$ e \qty{-9}{V}, um ciclo a cada
\qty{20}{ms}. O teste do relance: plana contra ondulada.
\end{solution}

\begin{solution}{exo:g9:alternating-current:8}
As máquinas girantes fazem CA naturalmente; e só a CA alimenta o
transformador, sem o qual a eletricidade não conseguiria atravessar o
país com \omterm{def:g8:voltage-voltmeter:voltage}{tensão} alta e entrar nas casas com uma \omterm{def:g8:voltage-voltmeter:voltage}{tensão} mansa.
\end{solution}

\begin{solution}{exo:g9:alternating-current:9}
Na pedalada lenta, os poucos ciclos por segundo do dínamo fazem a
lâmpada pulsar visivelmente --- cada crista um lampejo, cada zero uma
queda; pedalar mais rápido multiplica os pulsos até o olho fundi-los.
As $24$ imagens do cinema fixam a taxa aproximada de \omterm{def:g6:water-states:changes}{fusão} do olho:
abaixo de uns \qty{25}{Hz} de pulsos de \omterm{def:g1:light-and-shadows:source}{luz} a piscada incomoda; acima
disso, o brilho é lido como constante.
\end{solution}

\begin{solution}{exo:g9:alternating-current:10}
Cristas de cerca de \qty{325}{V} --- e cem delas por segundo,
cinquenta positivas e cinquenta negativas.
\end{solution}

\begin{solution}{exo:g9:alternating-current:11}
$T = 1/60 \approx \qty{0.017}{s}$; pico
$\approx 120 \times 1.4 \approx \qty{170}{V}$. Aquecedores e lâmpadas
se adaptam com indiferença; os carregadores modernos trazem
``100--240 V, 50/60 Hz'' nas etiquetas deles e convertem sem
problema --- os reclamantes são os aparelhos antigos de \omterm{def:g8:voltage-voltmeter:voltage}{tensão} única
e qualquer relógio motorizado que conte os ciclos da rede: ele
adiantaria com sessenta.
\end{solution}

\begin{solution}{exo:g9:alternating-current:12}
Brilhar mais nas duas cristas dá $2 \times 50 = 100$ pulsos por
segundo. Uma câmera de $25$ quadros amostra a cada \qty{0.04}{s} ---
quatro pulsos por quadro, mas não alinhados de modo idêntico em todo
ponto da exposição que rola: o pequeno descompasso entre o ritmo da
câmera e o da lâmpada pinta faixas que rastejam pelos quadros
sucessivos --- um batimento entre dois relógios. A \qty{60}{Hz}
($120$ pulsos), o descompasso contra os $25$ quadros muda, e as
barras mudam de largura e de \omterm{def:g7:motion-average-speed:formula}{velocidade} --- o padrão entrega a rede
local.
\end{solution}

\begin{solution}{pb:g9:alternating-current:1}
\textbf{1.} CC em $2.25 \times 2 = \qty{4.5}{V}$: a bateria chata de
bolso.
\textbf{2.} CA; $T = \qty{20}{ms}$, $f = \qty{50}{Hz}$; pico
$3 \times 2 = \qty{6}{V}$.
\textbf{3.} CC de \qty{4.5}{V} ligada \emph{ao contrário} --- as
pontas trocadas: inofensivo num \omterm{met:g9:alternating-current:scope}{osciloscópio}, instrutivo no
livro-caixa.
\textbf{4.} CA; $T = \qty{10}{ms}$, $f = \qty{100}{Hz}$; pico
\qty{3}{V}. Os pulsos dobrados de cem por segundo dela ficam além da
taxa de \omterm{def:g6:water-states:changes}{fusão} do olho com mais folga do que a piscada de cinquenta
ciclos de B.
\textbf{5.} O \omterm{def:g8:voltage-voltmeter:voltmeter}{voltímetro} informa o valor eficaz --- o equivalente
contínuo para o aquecimento ---, que, para uma senoide, fica em torno
do pico dividido por $1.4$: $6 \div 1.4 \approx \qty{4.3}{V}$. Os
dois números descrevem uma oscilação só.
\textbf{6.} Só a fonte A (ligada no sentido certo --- e a C depois de
plugada de novo). As fontes de CA precisam de um conversor --- o
bloco retificador do carregador --- entre elas e qualquer bateria.
\textbf{7.} As cinquenta oscilações por segundo da rede ficam perto
do próprio ritmo elétrico do coração e podem perturbá-lo: \omterm{def:g8:intensity-ammeter:intensity}{ampère} por
\omterm{def:g8:intensity-ammeter:intensity}{ampère}, a oscilação ameaça mais do que a marcha constante.
\textbf{8.} A senoide --- a oscilação perfeita, desenhada
naturalmente por uma bobina girando com \omterm{def:g7:motion-average-speed:formula}{velocidade} constante num
campo magnético: o alternador do próximo capítulo.
\textbf{9.} Um ciclo a cada $4$ quadradinhos (\qty{20}{ms}); cristas
de $325 \div 100 = 3.25$ quadradinhos acima e abaixo.
\textbf{10.} $f = 1/0.02 = \qty{50}{Hz}$; eficaz
$\approx 325 \div 1.4 \approx \qty{230}{V}$ --- a etiqueta,
recuperada da tela.
\textbf{11.} Uma linha plana em \qty{230}{V} --- a constância de
igual potência térmica.
\textbf{12.} Por exemplo: ``Uma linha plana é a constância de uma
\omterm{def:g3:first-electric-circuit:battery}{pilha} --- um sentido, um valor. Uma onda é a oscilação da rede ---
sentido e valor em ciclo incessante. E a taxa de câmbio justa entre
as duas é a \omterm{prop:g9:alternating-current:effective}{tensão eficaz}: \omterm{def:g5:heat-and-insulation:heat}{calor} igual, valor igual.''
\end{solution}
